\clearpage
\chapter*{第182次报告勘误}
\markboth{第182次报告勘误}{第182次报告勘误}
\addcontentsline{toc}{chapter}{第182次报告勘误}
\sid{e182}
\setcounter{footnote}{1}
\src{298}

\noindent\textbf{2. 形式几何与代数几何。}

\begin{center}
[布尔巴基讨论班，第11卷，1958/59年度，第182号，28页。]
\end{center}

\noindent\textbf{第8页，定理6。} --- 假设“$f$ 拟射影”可以换成较弱的假设
“$f$ 分离”；这是由下述结果（参见 SGA VIII，6.2）所得：每个拟有限且分离的
态射 $f:X\longrightarrow Y$ 都是拟射影的。

\src{299}
\noindent\textbf{第12页及第14页评注1。} --- 须指出，J.-P. Serre 构造了
代数闭域 $k$ 上的一个三维非奇异射影簇，该域的特征为 $p>0$；它并非由一个
整局部环上的固有概形约化而来，而该环的剩余域为 $k$、分式域的特征为零。
\footnote{Serre (Jean-Pierre). --- Exemples de variétés projectives en
caractéristique $p$ non relevables en caractéristique zéro, \emph{Proc. Nat.
Acad. Sc. U.S.A.}, t. 47, 1961, p. 108--109.}据说 Mumford 已找到一个以
非奇异射影曲面为例的类似结果。

\noindent\textbf{第15页，评注2和3。} --- 报告人目前对这里所猜想的结果已不
那么乐观。不过，评注3末尾提到的射影空间结构变形问题，已由 Hironaka 肯定地
解决；阿贝尔簇的相应问题则由 Koizumi 解决。

\noindent\textbf{第16页，第3行。} --- 须指出，Mumford 最近构造了亏格为
$g$ 的曲线的模概形（参见 SHMI）。此外，定理10证明，模理论中的“$n$ 级
Jacobi 概形”是非奇异的（甚至在 $\mathbf Z$ 上光滑）。

\noindent\textbf{第23页，第12行，推论5的第1行。} --- 补入：“$X$ 或 $Y$
在 $k$ 上固有”。

第1至5节的实质包含在 EGA III 已出版的部分，第6和7节的实质包含在
SGA III。关于基本群的研究，见 SGA V、IX、X、XI；关于 Lefschetz 型定理
及许多悬而未决的问题，另见 SGA 1962（报告 X、XII、XIII）。只有驯分歧
覆盖理论（参见定理14）尚未正式写成。定理14的推论给出了对如下对象的完全
分类：代数闭域上代数曲线的 Galois 覆盖，其阶与该域的特征互素。该推论曾在
三个场合发挥关键作用：

1）在 Igusa 对任意特征下非奇异射影曲面的 Picard 不等式的证明中；

2）在对下述群的研究中（Ogg 与 Šafarevič 各自独立开展了这一研究）：
任意特征下一元函数域上定义的阿贝尔簇作用下主齐性丛所成的群；

3）在 Artin 最近对代数簇“Weil 上同调”的若干关键定理的证明中。

\par
